% =============================================================================
% chapters/05_discussion.tex — Discussion (Feline Theme)
% =============================================================================

\section{Discussion}

This section interprets the findings of our study in the context of project management literature, highlights the theoretical and practical contributions of the FSM framework, and discusses the limitations of observing five uncooperative felines.

\subsection{Interpretation of Findings}

\subsubsection{Treat-Based Incentive Schemes}
Our results show a strong correlation between treat consumption ($T_c$) and purr volume ($d_{\text{purr}}$). This confirms that a transaction-oriented relationship exists between human project managers and feline stakeholders. However, the treat-to-nap relationship is complex: while salmon treats increase short-term purring, excessive treat delivery leads to immediate sleep induction, which halts project velocity. This represents a classic project management trade-off between resource utilization and productivity.

\subsubsection{Limitations of Laser Pointer Directing}
Although laser pointers are effective for short-term redirection of feline resources (H2 supported), qualitative logs suggest a high rate of stakeholder fatigue. By day 5, participants (especially Luna) realized that the red dot was uncatchable, leading to a significant decrease in responsiveness and an increase in aggressive ankle-swiping. This aligns with standard motivation theories indicating that unachievable goals lead to team frustration.

\subsection{Theoretical Contributions}
This study contributes to Project Management literature by applying empirical methods to the literal concept of "herding cats." We show that standard human-centric agile frameworks (such as Scrum) can be adapted to feline environments under the name of Agile Scratching, providing a new dimension to stakeholder engagement theories.

\subsection{Practical Implications}
For human project managers operating in multi-feline households, we offer four key recommendations:
\begin{enumerate}
    \item \textbf{Milestone-Based Treat Scheduling:} Reserve high-value treats (e.g., wet tuna) for critical project milestones, such as successful video calls without screen-walking.
    \item \textbf{Respect the Catnap Sprint:} Do not schedule any project activities between the hours of 10:00 AM and 4:00 PM, as stakeholder availability is near zero.
    \item \textbf{Mitigate Environmental Disruptions:} The vacuum cleaner should only be operated during off-site hours, as its activation triggers a total team collapse (H4 supported).
    \item \textbf{Keyboard Allocation Allowances:} Allow felines to sit on the laptop keyboard for at least 10 minutes per day. This serves as a vital team-building exercise and prevents sudden claw attacks.
\end{enumerate}

\subsection{Limitations}
Several limitations must be acknowledged. First, the sample size ($N = 5$) was extremely small and limited to a single household, reducing generalizability to multi-cat sanctuaries or feral environments. Second, the study suffered from high researcher bias, as the primary investigator frequently paused data logging to tell the participants how cute they were. Third, the long-term impact of catnip exposure remains unexamined, and its use may carry ethical concerns regarding stakeholder consent.
